% Generated draft: requires mathematical and source review.
\documentclass[11pt]{article}
\usepackage[margin=25mm]{geometry}
\usepackage{fontspec}
\setmainfont{Latin Modern Roman}
\usepackage{amsmath,amssymb,mathtools,unicode-math}
\setmathfont{Latin Modern Math}
\usepackage{enumitem,needspace,xurl,hyperref}
\setlength{\parindent}{0pt}
\setlength{\parskip}{5pt}
\setlength{\emergencystretch}{3em}
\newcommand{\N}{\mathbb{N}}
\newcommand{\Z}{\mathbb{Z}}
\newcommand{\Q}{\mathbb{Q}}
\newcommand{\R}{\mathbb{R}}
\newcommand{\C}{\mathbb{C}}
\newcommand{\F}{\mathbb{F}}
\newcommand{\A}{\mathbb{A}}
\newcommand{\PP}{\mathbb{P}}
\newcommand{\E}{\mathbb{E}}
\newcommand{\eps}{\varepsilon}
\newcommand{\dd}{\,\mathrm{d}}
\begin{document}
\section*{UnsolvedMath ID: 30001384}
\textbf{Source problem code:} catalogue entry
\subsection{Definitions and mathematical statement}
Let \(n\geq 1\) be an integer and let \(Y\) be a real-valued random variable. An \emph{\(n\)-dimensional version} of \(Y\) is a real random vector \(X=(X_1,\ldots,X_n)\) for which there exists a function
\[
\gamma:\R^n\setminus\{0\}\longrightarrow(0,\infty)
\]
such that
\[
\sum_{i=1}^n a_iX_i\ \stackrel{\mathrm d}{=}\ \gamma(a)Y
\qquad\text{for every }a=(a_1,\ldots,a_n)\in\R^n\setminus\{0\}.
\]
Here \(\stackrel{\mathrm d}{=}\) denotes equality in distribution. The function \(\gamma\) is called a \emph{standard}. The probabilistic problem is to characterize, for each law of \(Y\) and each \(n\), all joint laws of such vectors \(X\) and all standards \(\gamma\) that can occur. No independence of the coordinates is assumed.

For the geometric formulation, let \(K\subset\R^n\) be an origin-symmetric star body: a compact set star-shaped about the origin whose radial function
\[
\rho_K(u)=\max\{r\geq0:ru\in K\},\qquad u\in S^{n-1},
\]
is continuous and strictly positive, where \(S^{n-1}=\{u\in\R^n:\sum_{i=1}^n u_i^2=1\}\). Its Minkowski functional is
\[
\|x\|_K=\inf\{t>0:x\in tK\},\qquad x\in\R^n.
\]
In particular, \(\|0\|_K=0\). Convexity is not assumed, so this functional need not be a norm.

A function \(F:\R^n\to\C\) is positive definite if, for every integer \(m\geq1\), every \(x_1,\ldots,x_m\in\R^n\), and every \(c_1,\ldots,c_m\in\C\),
\[
\sum_{j,k=1}^m c_j\overline{c_k}\,F(x_j-x_k)\geq0,
\]
where the sum is required to be real and nonnegative. Define
\[
\Phi(K)=\bigl\{f\in C([0,\infty),\R):x\longmapsto f(\|x\|_K)\text{ is positive definite on }\R^n\bigr\}.
\]
The geometric problem is to characterize \(\Phi(K)\) for every origin-symmetric star body \(K\), and to characterize those \(K\) for which \(\Phi(K)\) is nontrivial. The source does not define ``nontrivial.'' Existence of a nonconstant member is a possible intended meaning, but requires confirmation; every nonnegative constant function belongs to \(\Phi(K)\).

\emph{Qualification concerning the asserted equivalence.} The supplied record describes the two problems as equivalent and reports that a standard is identified with a Minkowski functional. The displayed probabilistic definition itself imposes no symmetry or nondegeneracy hypothesis on \(Y\), and imposes no continuity, evenness, or homogeneity on \(\gamma\). In particular, when \(Y=0\) almost surely, \(X=0\) almost surely satisfies that definition for every positive \(\gamma\), including functions that are not Minkowski functionals. Moreover, a real radial positive-definite function is even, whereas the probabilistic formulation initially allows an arbitrary law of \(Y\). Matching probability laws to the geometric formulation also requires attention to normalization at the origin, typically \(f(0)=1\), which is not imposed in the definition of \(\Phi(K)\). Thus the source's asserted equivalence is retained with these unresolved qualifications, rather than asserted unconditionally for all the data allowed by its wording.
\subsection{Short English statement}
For each real random variable \(Y\), characterize the joint laws of random vectors whose every nonzero linear combination is a positive rescaling of \(Y\), together with all possible rescaling functions. The source also asks for the continuous real functions of a star body's Minkowski functional that are positive definite, and for the bodies admitting a nontrivial such class; the meaning of nontriviality and the hypotheses connecting the two formulations require review.
\subsection{Sources}
\begin{itemize}
\item Supplied local UnsolvedMath registry record, permanent ID \(30001384\), title \emph{Positive-Definite Functions of Minkowski Functionals}, exact source problem code \texttt{catalogue entry}, in the \emph{Oberwolfach Reports Open Problems} collection. The record gives proposed year \(2009\) and does not identify an individual proposer.
\item Supplied report citation: \emph{Convex Geometry and its Applications (2010)}, Oberwolfach Reports, pp.~2851--2926, DOI \texttt{10.4171/owr/2009/53}. \url{https://doi.org/10.4171/owr/2009/53}. The literature-source label is ``original Oberwolfach report 4136.'' The supplied background describes these reports as documenting mathematical workshops and their open-problem discussions.
\item Supplied literature-source label: ``Koldobsky standards theorem.'' \url{https://arxiv.org/abs/0903.1433}.
\item Supplied literature-source label: ``Schoenberg positive-definite-function result.'' \url{https://doi.org/10.1112/S0024609399006037}.
\item Supplied dated literature curation, checked \texttt{2026-08-21}: the registry labels the problem \emph{partially solved} and states that Koldobsky proved Lisitsky's embedding-in-\(L_0\) conjecture, providing a strong necessary condition for nontrivial classes, while the full characterization of all standards and all classes \(\Phi(K)\) remains incomplete. This is the supplied dated assessment, not independently verified literature status or a verified theorem claim in this draft. The background explicitly requests independent verification of time-sensitive status and cited claims.
\item The supplied registry labels the statement \texttt{corrected\_verified} and reports that the official report defines an \(n\)-dimensional version, identifies \(\gamma\) as a Minkowski functional, and states the positive-definite-function formulation. These are supplied source-verification claims; no source was visited or independently checked in preparing this draft.
\end{itemize}
\end{document}
